\section{Scalar acquisition as a supporting specialization}\label{app:scalar}

Choose $0<h\le1$, a nonnegative $C^2$ cutoff $w$, supported in
$[-1,1]$ and equal to one on $[-1/2,1/2]$, with $w\le1$.
Choose $0\le\kappa\le1$, equal to one on $[0,1]$ and zero on
$[2,\infty)$, with bounded first and second derivatives. Define
\begin{equation}\label{eq:scalar-tests}
 a(x,y)=w((x-X)/h)y_+^2\kappa(y),\qquad
 \varphi_Q=2a,\qquad
 \varphi_\eta=2\log(1+\eta a)\quad(\eta\ge0).
\end{equation}
These compact $C^{1,1}$ tests are supported in
$|x-X|\le1$, $0\le y\le2$. The support lies strictly within
$D(c,r)$ since $\sqrt{17}<25/6$. At $y=0$ the value and the first
normal derivative vanish, so their weak Laplacians contain no
boundary delta.

The canonical strip, conjugation, and \eqref{eq:zero-measure} imply
\begin{equation}\label{eq:scalar-fold}
 Q_{X,h}=\nu_F(\varphi_Q),\qquad
 \log Z_{X,h}(\eta)=\nu_F(\varphi_\eta).
\end{equation}
The factor two accounts for lower conjugate copies. Real labels
contribute zero. The harmonic logarithm of $D$ disappears.
Define the comparison values
$\widehat Q=\nu_{\widehat F}(\varphi_Q)$ and
$\widehat L_\eta=\nu_{\widehat F}(\varphi_\eta)$.
The approximating zeros in $1<y<2$ keep their actual cutoff weights:
there is no strip theorem for the approximant.

\begin{theorem}[Complete one-point acquisition]\label{thm:scalar}
For $X\ge10$, $0<h\le1$, and the tests \eqref{eq:scalar-tests},
choose $J,U$ so that \eqref{eq:complete-delta} is at most
$\min(1,c_0/2)$. Then
\begin{equation}\label{eq:scalar-transfer}
 |Q_{X,h}-\widehat Q|\le K_QE_X,\qquad
 |\log Z_{X,h}(\eta)-\widehat L_\eta|\le K_\eta E_X,
 \qquad
 K_\varphi=\frac{\norm{\Delta\varphi}_\infty}{2\pi}.
\end{equation}
All multiplicities are included. The constants are explicitly
given by \eqref{eq:radii}, \eqref{eq:canonical-constants} and
\eqref{eq:complete-delta}.
\end{theorem}
\begin{proof}
The anchor and disk bounds establish Theorem~\ref{thm:log};
apply Corollary~\ref{cor:weak-transfer} and \eqref{eq:scalar-fold}.
\end{proof}

For $v(y)=y_+^2\kappa(y)$ the derivative costs are directly available:
\begin{align}
 \Delta a&=h^{-2}w''v+wv'',&
 |\nabla a|^2&=h^{-2}|w'|^2v^2+w^2|v'|^2,\label{eq:scalar-derivatives}\\
 \Delta\varphi_\eta&=
 2\left(\frac{\eta\Delta a}{1+\eta a}
              -\frac{\eta^2|\nabla a|^2}{(1+\eta a)^2}\right).
 \notag
\end{align}
In particular $K_Q\le(h^{-2}\norm{w''}_\infty\norm v_\infty+
\norm w_\infty\norm{v''}_\infty)/\pi$.
The denominators in \eqref{eq:scalar-derivatives} are at least one;
growing $\eta$ still incurs its displayed derivative cost.
The Jensen count also gives the coarse canonical bound
$N_w\le2n_X$, hence $Q_w\le2n_X$, because every upper label in
the window lies in $D(c,\rho_*)$, every lower nonreal label has an
upper partner, and real labels need be counted only once.

For a nonzero test, any desired statistic accuracy $e>0$ is achieved
by setting
\begin{equation}\label{eq:precision-choice}
 \delta_*=\min\left\{1,\frac{c_0}2,\
    M_X\left(\frac{e}{2K_\varphi C_X}\right)^{1/\alpha_X}\right\}
\end{equation}
and choosing finite cutoffs satisfying
\begin{equation}\label{eq:tail-allocation}
 5C_2(J)I_{b_*}\le\frac{\delta_*}{2\mathcal B_X},
 \qquad B_0I_{b_*}^{>U}\le\frac{\delta_*}{2\mathcal B_X}.
\end{equation}
Then $K_\varphi E_X\le e/2$. This proves finite precision exists,
without claiming a uniform computation time. If an evaluator returns
$\widehat q$ with $|\widehat q-\widehat Q|\le E_{\rm eval}$,
the actual exclusion condition for a plateau pair of height at least
$d$ is
\begin{equation}\label{eq:scalar-consumer}
 \widehat q+E_{\rm eval}+K_QE_X<2d^2 .
\end{equation}
There is an analogous condition with margin
$2\log(1+\eta d^2)$. Quadrature and rounding belong in
$E_{\rm eval}$. No favorable $\widehat q$ is obtained from the
approximation theorem itself.
